\documentclass[12pt, twocolumn, landscape, a4paper]{article}

\usepackage[margin=1in]{geometry}
\usepackage[utf8]{inputenc}
\usepackage[spanish,es-noshorthands]{babel}
\usepackage{amssymb, amsmath, amsbsy}
\usepackage{mathpazo}
\usepackage{graphicx}
\usepackage{tikz}
\usetikzlibrary{angles,quotes}
\usepackage[pdftex]{hyperref}
\hypersetup{pdftitle={Producto Escalar},pdfauthor={Luis Carrillo Gutiérrez}}

\renewcommand{\familydefault}{\sfdefault}
\pagestyle{empty} 

\begin{document}
\subsubsection*{Producto Escalar entre vectores}
\begin{itemize}
\item{Si \scalebox{1.3}{$\overrightarrow{m}$} y \scalebox{1.3}{$\overrightarrow{n}$} son \textbf{perpendiculares} ($\theta = 90^{\,\circ}$) \\ 

\begin{tikzpicture}
	\coordinate (A) at (0, 0);
	\coordinate (B) at (-3, 2);
	\coordinate (C) at (1.86, 2.5);
	
	\begin{scope}[-latex]
	\draw (A) -- (B);
	\draw (A) -- (C);
	\end{scope}
	
	\draw pic[draw,angle radius=.35cm] {right angle=C--A--B};
	\draw (-2.15, 1) node {$\vec{m}$};
	\draw (1.25, 1.1) node {$\vec{n}$};
	\draw (0, .85) node {$90^{\,\circ}$};
\end{tikzpicture}

$\vec{m} \cdot \vec{n} = m \cdot n \cdot \cos 90^{\,\circ} = 0$
}
\item{Si $\overrightarrow{m}$ y $\overrightarrow{n}$ son \textbf{paralelos} y con el \textbf{mismo sentido} ($\alpha = 0^{\,\circ}$) \\ $\vec{m} \cdot \vec{n} = m \cdot n \cdot \cos 0^{\,\circ} = a \cdot b$}
\item{Si $\overrightarrow{m}$ y $\overrightarrow{n}$ son \textbf{paralelos} y con \textbf{distinto sentido} ($\alpha = 180^{\,\circ}$) \\ $\vec{m} \cdot \vec{n} = m \cdot n \cdot \cos 180^{\,\circ} = -1 \cdot a \cdot b $}
\end{itemize}
\end{document}
